\cotefile{161-4}
\scanpage{3}
\mspage{1}\mspage{\textcircled{\unsure{1}}}

\uline{Descente plate}

Module plat sur un anneau \deleted{\unsure{ou fid. plat}}

\margin{\uline{Faisceautisation} \inserted{de la notion de platitude} \uline{sur les espaces} \unsure{s/schéma|i.e. pour} $f:X\mathrel{-}Y$.}

Algèbre plate $B/A$

\uline{Proposition} $B$ algèbre sur $A$. Alors $B$ fid. plat sss $B$ est \uline{plat} et $\operatorname{Spec}B\mathrel{-}\operatorname{Spec}A$ est \uline{surjectif} \deleted{\unsure{(En fait, le}} [ou encore, $\forall$ idéal [premier] $\unsure{I|J|\mathfrak J}$ de $A$, $A/\unsure{IA|JA|\mathfrak JA}\longleftarrow B/\unsure{IB|JB|\mathfrak JB}$

\deleted{\uline{Néc} \unsure{Car} $A/p\longleftarrow$}

$B$ fid plat/$A$ \unsure{En effet}

$\Rightarrow\quad A\longrightarrow B$ injectif (car \deleted{le noyau $I$ doit satisfaire} \inserted{si $I\subset A$ est un idéal, non nul} \unsure{$I\mathrel{-} A$ inj.} \deleted{\unsure{$I=0$}} $\underset{\ne0}{I\otimes_A B}\mathrel{-}B$ injectif donc $IB\ne0$)

% margin note written sideways
\margin{$M$ $A$-module. $M$ fid plat $\Longleftrightarrow$ $\forall p$ premier \deleted{\unsure{de}} $A$ \inserted{(ou \unsure{seult} maximal)} $M\otimes_A k(p)\ne0$ néc. \uline{Suffisant} : \ldots\ \uline{Cor} Cas $B$ une $A$-algèbre. $B$ fid plat $\Longleftrightarrow B$ plat et $\operatorname{Spec}B\mathrel{-}\operatorname{Spec}A$ surj. \uline{Cor} $A\mathrel{-}B$ homom. local \unsure{(locaux)} $B$ plat $\Longleftrightarrow B$ fid plat \unsure{(voyez modules, ou} \unsure{$M$ $\mathfrak m$-module de t. f.)}.}

\[
\forall\ \text{idéal } I\subset A,\qquad I=IB\cap A \qquad (A/I\longrightarrow B/IB\ \text{injectif})
\]

p. ex. si $I$ premier \quad $A/p\deleted{\unsure{=}}\longleftarrow B/pB$ \quad \uline{injectif}

\deleted{Donc \unsure{il} $\operatorname{Spec}B/pB$} \inserted{\unsure{i.e.|Ind.}} $\exists$ \uline{un idéal premier} $q$ \uline{de} $B$ \uline{qui induit} $p$

donc $\operatorname{Spec}B\mathrel{-}\operatorname{Spec}A$ surjectif

\unsure{Inv.|Inversement} si $\operatorname{Spec}B\longrightarrow\operatorname{Spec}A$ surj. et $B$ plat sur $A$, prouvons qu'il est \uline{fid} \uline{plat} : $M$ un $A$-module $\ne0$, prouvons $M\otimes_A B\ne0$.

\deleted{\unsure{Soit supposons|On peut supposer} $M$ de t.f. (} Soit $M'\subset M$ \deleted{de t.f. $\ne0$}, monogène \deleted{de t.f.} \inserted{$\ne0$}, il suffit de prouver $M'\otimes B=0$ (car il est $\subset M\otimes B$). Il a un quotient de la forme $A/\mathfrak m$ (un idéal maximal), alors \unsure{on} \deleted{peut supposer $M$} est ramené : prouvons que $B/\mathfrak mB\ne0$, clair car $\mathfrak m$ est induit par un idéal premier de $B$

\uline{Déf} Soit $f:X\mathrel{-}Y$ morphisme de \deleted{\unsure{schémas|espaces annelés}} \deleted{\unsure{préschémas|schémas}} schémas. \deleted{$M$ \unsure{un} $\underline O_X$} On dit que \deleted{$M$ plat/$Y$ sss} $f$ est fid plat sss \quad /

a) b) pour morphismes plats et fid plats

\uline{Prop} Si $f:X\longrightarrow Y$ est plat, \inserted{a)} il induit \deleted{des} \deleted{\unsure{une}} surjection $\operatorname{Spec}\underline O_{X,x}\longrightarrow\operatorname{Spec}\underline O_{Y,y}$ \quad b) Tout pt maximal de $X$ est au dessus d'un pt \uline{maximal} de $Y$

% archive sheet number: 2
\scanpage{4}
\mspage{2}

\uline{Cor} Si $f$ est loc. de présentation finie et plat, $f$ est \inserted{\uline{univ}} \uline{ouvert}.

\uline{Prop} Si $S'\overset{f}{\longrightarrow}S$ est fid plat, l'application $T\longmapsto f^{-1}(T)$ des sous-schémas de $S$ vers les sous-schémas de $S'$ \deleted{est \unsure{injectif}} est injective.

\uline{\unsure{Cor|Prop}} Si $X,Y$ sur $S$, alors $\operatorname{Hom}_S(X,Y)\mathrel{-}\operatorname{Hom}_{S'}(X',Y')$ injective.

\deleted{Lemme}

car.

\diagram{\begin{tikzcd}[ampersand replacement=\&]
X \arrow[r,dash] \arrow[d,"f"'] \& X' \arrow[d,"f'"]\\
Y \& Y' \arrow[l,"\text{\uuline{plat}}"]
\end{tikzcd}}

\[
 f_*(F)'\overset{\sim}{\longrightarrow}f'_*(F')
\]

\diagram{\begin{tikzcd}[ampersand replacement=\&]
Z' \& Z \arrow[l] \\
S' \arrow[r,"f"'] \& S
\arrow[from=1-1,to=2-1,"g'"'] \arrow[from=1-2,to=2-2,"g\ \text{quasi-compact}"]
\end{tikzcd}}

$g(Z)=T$

\uline{Th} Soit $f:S'\longrightarrow S$ un morphisme plat. Pour tte partie \deleted{\unsure{la}} \deleted{\unsure{loc. fermée}} $T$ de $S$ \deleted{telle que} \inserted{image d'un morphisme} $Z\longrightarrow S$ quasi-compact, on a $f^{-1}(\overline T)=\overline{f^{-1}(T)}$.

\margin{\uline{Dém} \unsure{OPS|SPS} $X,Y,Z$ affines.}

\diagram{\begin{tikzcd}[ampersand replacement=\&]
C' \& C \arrow[l]\\
\unsure{I|J}\quad A' \arrow[u] \& A\quad I \arrow[l] \arrow[u,"u"]
\end{tikzcd}}

\margin{Alors $\overline T=V(I)$ ($I=\operatorname{Ker}u$) $f^{-1}(\overline T)=V(IA')$ \quad $\overline{f^{-1}(T)}=V(K)$ \quad $K=\operatorname{Ker}u$ \quad $\overline{g'(Z')}$}

\margin{\uline{Autre dém} On est ramené au \uline{cas} $g$ dominant $\Rightarrow g'$ dominant. \unsure{/ et|et} $\Leftarrow$ $g$ dominant, $S'\ne\varnothing\Longrightarrow Z'\ne\varnothing$. Mais cela est vrai pour \deleted{tt} morphismes \uline{générisants} car on est ramené aussitôt aux \uline{cor. Univ.} Soit $g:Z\longrightarrow S$ qu. cpt. Alors dominant \deleted{\unsure{sss|Mais}} pour tt pt maximal \deleted{\unsure{$s$}} $\unsure{\in S|\text{de }S}$ $Z_s\ne\varnothing$ \deleted{\unsure{est}}.}

\uline{Cor} Si $f$ fid plat \inserted{et quasi-compact}, la topologie de \deleted{$S$} \inserted{$f(S')$ induite par $S$} est quotient de celle de $S'$.

\deleted{$\unsure{T=f^{-1}(f(T))}\quad\overline T\quad f^{-1}(\overline{f(T)})\quad\unsure{\Longleftrightarrow}\quad\unsure{\text{si }}T=\overline T,\quad T\unsure{=}f^{-1}(f(T))$}

Il faut prouver que si $T'$ est une fermée de $S'$ saturée pour $f$, alors $T=f(Z')$ est rel. fermé dans $Z=f(S')$, i.e. que $\overline T\cap Z=T$ i.e. $\underset{\Vert\,T'}{f^{-1}(T)}=\underset{\Vert\,\overline{T'}}{f^{-1}(\overline T)}$ O.K.

\inserted{Si $f:S'\mathrel{-}S$ fid plat, qu. cpt \inserted{(ou mieux, gén., surjectif)}}

\uline{Cor}

\[
\begin{aligned}
 F(S)&\mathrel{-}F(S')= F(S'')\\
 O(S)&\mathrel{-}O(S')= O(S'')\\
 LF(S)&\mathrel{-}LF(S')= LF(S'')\\
 [LF^{\unsure{rc|qc}}(S)&\longrightarrow LF^{rc}= LF^{rc}(S'')\ ?\,]
\end{aligned}
\]

(où $F={}$ parties fermées \quad $O:{}$ parties ouvertes \quad $LF$ parties loc. fermées \inserted{rétrocompactes}) est \uline{exact} \quad $S''=S'\times_S S'$

% slanted
\margin{en fait aussi pour les parties \unsure{``constructibles''|structure induite} \deleted{\illegible{1}}}

\unsure{Bien Etd|Bien Ets} si $f$ surj, on a \quad $P(S)\mathrel{-}P(S')= P(S'')$ \quad ($P={}$ parties) car $S''\mathrel{-}|S'|\times_{|S|}|S'|$ surjectif. Donc l'exactitude des \deleted{deux} \inserted{deux} premiers exprime simplement que la top de $S$ est quotient, \deleted{\unsure{le}} \unsure{car|cas} est \unsure{bijectif|trivial|injective}.

\uline{Cor} Une partie \inserted{$T$} de $S$ est \unsure{ouverte|fermée} (fermée, loc fermée rétrocompacte) sss $f^{-1}(T)$ l'est. Si $T$ est l'image d'un morphisme qu. cpt, alors $T$ est loc. f. sss $f^{-1}(T)$ l'est\ldots

Pour le troisième cas, on est ramené aux deux premiers en introduisant l'adhérence et réutilisant le Th.

% archive sheet number: 3
\scanpage{5}
\deleted{Descente} \mspage{3}

% slanted
\margin{\textcircled{1} \uline{descente de propriétés de finitude}}

\uline{Prop} \uline{Soit} $S'\overset{f}{\mathrel{-}}S$ \uline{surjectif}. Pour qu'un $S$-morphisme $g:X\mathrel{-}Y$ soit surjectif (resp. radiciel) il f et s que $g':X'\mathrel{-}Y'$ le soit. [Si $f'$ inj $\mathrel{-}$ bij $\mathrel{-}$

\uline{Prop} \deleted{Soit} $f:S'\mathrel{-}S$ surjectif qu. cpt \inserted{(resp. fid plat qu. cpt)}. Alors $f$ est quasi-compact \inserted{ou quasi-séparé, i.e. $\Delta_f$} (resp. de t.f. \unsure{)}, loc. de t.f., de prés finie, loc de prés. finie) sss $f'$ l'est

\diagram{\begin{tikzcd}[ampersand replacement=\&,column sep=small]
\& X' \arrow[d] \& \& X \arrow[d,dash]\\
\unsure{Y'} \arrow[r,dash]\arrow[dr,dash] \& Y' \arrow[dr,dash] \& \unsure{X'} \arrow[r,dash]\arrow[dl,dash] \& X \arrow[r,dash]\arrow[d,dash] \& Y\arrow[dl,dash]\\
\& S'\arrow[rr,dash] \& \& S
\end{tikzcd}}

La partie non \unsure{resp.|respée} immédiate \deleted{\illegible{1}} On est ramené pour la partie \unsure{resp.|respée} : \unsure{au}.

\uline{Lemme} Soit $A\mathrel{-}A'$ fid plate, $B$ $A$-algèbre. Alors $B$ est de t.f. \inserted{(p.f.)} sss $B'$ l'est.

Les \unsure{$\Rightarrow$|mêmes} arguments, par modules, \unsure{donc|donnent}

\uline{Cor.} Soit $f:S'\mathrel{-}S$ fid plat qu. cpt, $\unsure{\mathcal F|F}$ faisceau qu. ch. sur $S$. Alors $\mathcal F$ de t.f. (pr. f.) sss $\mathcal F'$ l'est.

On prouve $\unsure{=|\text{même}}$ chose pour \unsure{libre|loc.}. libre de t.f. \inserted{loc libre de rang $n$ \quad donc : inversible}

\margin{\uline{Prop} $f$ fid plat. Alors $g$ est plat (fid plat) sss $g'$ l'est \quad ou \unsure{par modules|on peut prendre modules}}

% slanted
\margin{\textcircled{2} \uline{descente de propriétés topologiques} et conséquences\ldots}

Th.

\uline{Cor} \inserted{Soit $S'\mathrel{-}S$ fid plat \inserted{(ou univ. gén.)} et quasi-compact} Soit $f:X\longrightarrow Y$ un \unsure{morphisme|S-morphisme}. Alors \deleted{\unsure{Alors}} \inserted{\unsure{Si|$S'$}} $f'$ est ouvert (resp \inserted{générisant}, fermé, resp. quasi-compact et un homéo dans, resp. un homéo), $f$ de m.

Donc $f$ est imm ouvert (resp. \unsure{--}) sss $f'$ l'est.

\uline{Cor} $f$ séparé (propre) sss $f'$ l'est

\diagram{\begin{tikzcd}[ampersand replacement=\&]
 X'\arrow[r]\arrow[d,"g'"'] \& X\quad U\arrow[d,"g\ \text{surj}"]\\
 Y'\arrow[r,"f_g\ \text{surj}"'] \& Y
\end{tikzcd}}

\diagram{\begin{tikzcd}[ampersand replacement=\&]S'\arrow[r,dash,"{\text{surj}}","f"'] \& S\end{tikzcd}}

% slanted
\margin{\uline{Quelques} \uline{autres propriétés}}

\uline{Cor} $f$ un iso (mono) sss $f'$ l'est.

On est ramené au cas d'un iso, \deleted{et pour cela on est ramené} \inserted{On sait déjà par ce qui \unsure{pr.}} \deleted{\unsure{au fait} que si $f'$ l'est, $f$ est \unsure{a priori} ($f$ l'étant), donc} est un homéomorphisme, reste à prouver que $O_Y\mathrel{-}f_*(O_X)$ est un iso, or cela le devient après chgt. de base, o.k.

\deleted{\unsure{Cor} Soit $f$ fid plat, \illegible{1}}

\deleted{\uline{Question} Si $g'$ une imm, $g$ une ? O.K. si $g$ \unsure{loc. de prés. finie} ?}

surj, radiciel, \deleted{\unsure{quasi-sép}} \deleted{\unsure{injectif}} \inserted{plat, fid plat} qu cpt, qu. sép, \inserted{\unsure{cohér|coh.}}, t.f., p.f. loc t.f., loc p.f. séparé, propre, \deleted{\illegible{1}} univ. fermé, univ. ouvert, homéo univ, \unsure{domin.} biration\unsure{nel}

\unsure{Diagonal} \unsure{imm|immersions} (imm. fermée, imm. ouverte)

\deleted{\illegible{1}} affines, entiers, finis, quasi-affines, quasi-finis \deleted{plat fid}

[lisse, net, étale formellement lisse (?), formellemt étale\ldots.

\uline{TSVP}

% archive sheet number: 4
\scanpage{6}
\mspage{4}

\textcircled{3} \uline{Descente de} \deleted{\unsure{ses}} \uline{faisceaux}, \deleted{\unsure{d'immersions}} \deleted{\unsure{de morphismes}} \uline{qu. coh.}, \deleted{\uline{et de morphismes}} \inserted{de \deleted{\unsure{morphismes}} sous-schémas}, de morphismes de schémas, de schémas affines.

\deleted{\uline{Prop} $f:S'\mathrel{-}S$ fid plat et qu. cpt, $F$ qu. coh. sur $S$. Alors}

\[
\deleted{\operatorname{Ssf}(F)\mathrel{-}\operatorname{Ssf}(F')\rightrightarrows\operatorname{Ssf}(S'')\quad\text{est exact.}}
\]

\deleted{On est ramené au cas $S$ affine, puis $S'$ affine, et au résultat \unsure{connu|anc.|conn.} d'algèbre commutative.}

\diagram{\begin{tikzcd}[ampersand replacement=\&,row sep=small]
N'\cap M=N_0 \& N' \& {}\\
 M \& M'\arrow[u] \& M''\\
 A\arrow[r,dash] \& A'\arrow[r,dash]\arrow[r,shift right=.5ex,dash] \& A''
\end{tikzcd}}

\uline{Th} Soit $f:S'\mathrel{-}S$ fid plat qu. cpt. Alors $f$ est de descente effective (univ.) pour la catégorie fibrée des Modules qu. coh.

Expliquer ce que cela signifie.

\uline{Cor} \deleted{$f$ est de \unsure{d.}} Soit $F$ qu. coh sur $S$. Alors suite exacte

\[
\operatorname{Ssf}(F)\mathrel{-}\operatorname{Ssf}(F')=\operatorname{Ssf}(F'')
\]

\uline{Corollaire} On a des suites exactes

\deleted{\unsure{``ssf$(S)$''|$\operatorname{Ssf}(S)$}}

\[
\begin{aligned}
 \operatorname{Ss}\deleted{\unsure{F}}f(S)&\mathrel{-}\operatorname{Ssf}(S')\rightrightarrows\operatorname{Ssf}(S'')\\
 \operatorname{Ss\,ou}(S)&\longrightarrow\operatorname{Ss\,o}(S')\rightrightarrows\operatorname{Ss\,o}(S'')\\
 \operatorname{Ss\,rc}(S)&\mathrel{-}\operatorname{Ss\,rc}(S')\rightrightarrows\operatorname{Ss\,rc}(S'')
\end{aligned}
\]

\margin{Peut-on descendre les sous-schémas \unsure{par imm. rétroc.|pas nécess. rétroc.} \unsure{ou ouvertes|ni ouvertes|non ouverts} ?}

\uline{Corollaire} $f$ est de descente (\unsure{univ|univ.}) pour la \unsure{cat-cart|catégorie} fibrée des schémas sur d'autres :

\[
 \operatorname{Hom}_S(X,Y)\mathrel{-}\operatorname{Hom}_S(X',Y')=\operatorname{Hom}_S(X'',Y'')
\]

\diagram{\begin{tikzcd}[ampersand replacement=\&]
 X\times_S Y=Z\arrow[d,dash] \& Z'\arrow[d,dash] \& Z''\arrow[d,dash]\\
 Y\arrow[u,bend left] \& Y' \& Y''\\
 S\arrow[r,dash] \& S'\arrow[r,shift left=.65ex,dash]\arrow[r,shift right=.65ex,dash] \& S''
\end{tikzcd}}

On le sait par le corollaire précédent si $Y$ est qu. sép sur $S$. Mais on se ramène aisément au cas $S,Y$ affines\ldots

On retrouve que $g'$ un iso $\Rightarrow g$ un iso

\uline{Corollaire} Pour que $f'$ une imm. fermée (ouverte, \inserted{arbitraire} \deleted{\unsure{propre}}), il f et \uline{s que} $f$ le soit

\deleted{\unsure{Ces prop gén si $g'$ une imm fermée et $g$} \unsure{aussi ? O.K. si $g$ un morphisme}}

\uline{Corollaire} \deleted{des} d'Algèbres qu. coh.

$\Downarrow$ \deleted{Cor} \uline{Th} \inserted{Si $f$ fid plat qu. cpt, alors} \deleted{\unsure{$f$ est fid}} un morphisme de descente eff. (univ) pour les schémas relatifs \uline{affines}

\uline{Cor} $g$ \deleted{une immersion} affine $\Longleftrightarrow g'$ affine

\[
\begin{aligned}g\ \text{quasi-affine}&\Longleftrightarrow g'\ \text{quasi-affine}\\
 g\ \text{entier (fini, \inserted{fini loc. libre})}&\Longleftrightarrow g'\ \text{entier (fini, fini loc. libre)}
\end{aligned}
\]

\margin{[car $g$ qu. aff $\Longleftrightarrow g$ \unsure{qch|qu. cpt|coh.} et $X\mathrel{-}\operatorname{Spec}g_*(O_X)$ une imm}

% archive sheet number: 5
\scanpage{7}
\mspage{5}

Pour ce dernier point, on est ramené au

\uline{Lemme} $A'$ fid plat sur $A$, $B$ une $A$-algèbre. Alors $B$ fini (entier) sur $A$ sss $B'$ l'est sur $A'$.

Fini est une histoire de module. Entier s'y ramène en regardant les $A[x]$, $x\in B$.

\textcircled{4} \uline{Descente de l'amplitude}

\diagram{\begin{tikzcd}[ampersand replacement=\&]
\underline L\quad X\arrow[d,dash] \& X'\arrow[d,dash]\\
Y\arrow[r,dash]\arrow[d,dash] \& \unsure{Y''|Y'}\\
S \& S'\arrow[l,"{\substack{f\ \text{fid plat}\\\text{qu cpt}}}"']
\end{tikzcd}}

\uline{$L$} (très) ample rel /$Y$ $\Longleftrightarrow$ $L$ \inserted{(très)} ample rel. à $Y$

On peut supposer $X$ cohérent sur $Y$, alors $g_*(L)$ \unsure{qu. cpt|qu. coh} et $g_*(L)'\simeq g'_*(L')$

Très ample signifie que 1) $g^*(g_*(L))\longrightarrow L$ épi, donc $X\mathrel{-}P(g_*(L))$ défini. \deleted{\unsure{2)}} 2) \deleted{$g^*(g_*(L))\unsure{\mathrel{-}L}$} $X\mathrel{-}P(g_*(L))$ une imm.

Si $g$ de t.f., ``ample'' se ramène : ``très ample'', donc c'est.

Pour le cas général, \inserted{on peut supposer $Y$ qu. cpt,} on utilise le critère \unsure{que|avec} pour $F$ qu. ch. de t.f. sur $S$. et $n$ grand, $f^*(f_*(F(n)))\mathrel{-}F(n)$ \unsure{épij.|épi}

On retrouve que $g$ quasi-affine $\Longleftrightarrow g'$ qu.-affine

\uline{NB} les seuls exemples de propriétés de morphismes qui ne se descendent pas par fpqc sont les morphismes \uline{projectifs} et \uline{quasi-projectifs}

\textcircled{5} \uline{Questions} d'\uline{effectivité} \uline{des données de descente}

% archive sheet number: 6
\scanpage{8}
\uline{Morphismes radiciels} \quad Soit $f:X\mathrel{-}Y$ un morphisme de schémas. Conditions équivalentes

a) $f$ univ. \uline{injectif}

b) $\forall y\in Y$, $X_y$ est radiciel sur $k(y)$, i.e. \deleted{\unsure{est}} a au plus un point $x$, et $k(x)$ est une ext. radicielle de $k(y)$

\unsure{c|b}) $\forall$ corps $K$, $X(K)\longrightarrow Y(K)$ est \uline{injective}

d) $\delta_f:X\longrightarrow X\times_Y X$ est surjective

Les morphismes sont dits \uline{radiciels}. (Un \uline{monom.} est radiciel, donc une imm, telle $\delta_f$, l'est) Satisfont à b) et c] ($gf$ radiciel $\Rightarrow f$ radiciel)

\uline{Prop} Soit $f:X\mathrel{-}Y$ un morphisme entier radiciel \inserted{surjectif}. Alors c'est un \uline{homéom.} \deleted{\unsure{univ}} \uline{universel}.

On peut prouver la réciproque -- mais c'est nettement plus délicat

\uline{Exemples} 1) $k'/k$ ext. radicielle (p. ex. clôture parfaite)

$X_{k'}\mathrel{-}X_k$ est entier radiciel surjectif, donc un homéom (universel)

2) Courbe à cusp, et sa normalisée.

\uline{Morphismes univ. ouverts}

\uline{Condition néc.} Pour qu'un morphisme soit ouvert : il \uline{est générisant}

Cette condition est \uline{suffisante} si $f$ loc. de prés. finie

Donc morphismes univ. ouverts \deleted{sont} \uline{univ. générisants}, et réciproqt si $f$ est loc. de prés finie

Morphismes univ. ouverts et univ. générisants satisfont 1), 2).

% archive sheet number: 7
\scanpage{9}
\uline{Schémas quasi-affines} \inserted{(\unsure{quasicompacts|quasi-compacts}, et isom. à sous-schéma ouvert d'un schéma affine)} / \uline{morphismes quasi-affines}

Conditions équivalentes sur schéma $X$ cohérent, $A=\Gamma(X,\underline O_X)$

a) $X$ quasi-affine

b) $X\overset{u}{\mathrel{-}}\operatorname{Spec}A$ une imm. ouverte

b') $X\mathrel{-}\operatorname{Spec}A$ un homéomorphisme dans

c) $\underline O_X$ est ample

\[
\left\{\begin{aligned}
 &\text{c') Les }X_f\text{ forment une base top. de }X\\
 &\text{c'') les }X_f\text{ affines forment base top. de }X\\
 &\text{c''') Tout } F\text{ quasi-cohérent \uline{engendré} par ses sections}\\
 &c^{\mathrm{iv})})\ \text{Id pour }F\text{ qu. ch }\subset\underline O_X
\end{aligned}\right.
\]

d) $\exists$ Faisceau \deleted{ample} \inserted{inv.} $L$ tel que $L$ et $L^{-1}$ amples.

d') Tout Module inv. $L$ est ample.

\uline{Morphismes} \uline{quasi-affines} : Stable par changemt de base, par composition les imm. quasi-compactes sont qu. affines, donc si $g$ qu. sép, $\delta_g$ qu. affine

$gf$ quasi-affine et $g$ qu. sép $\Longrightarrow f$ quasi-affine

\uline{Th} Si $f$ de type fini, \deleted{\unsure{séparé}} à fibres finies ($f$ quasi-fini) et $f$ séparé, donc $f$ quasi-affine (donc quasi-projectif)

\uline{morph. quasi-projectif} : de t.f et $\exists L$ $f$-ample.

Propre $={}$ univers. fermé, séparé, de t.f.

entier

\[
\begin{aligned}\text{projectif}&=\text{propre}+\text{quasi-projectif}\\
 &=\deleted{\unsure{\text{propre}}}\ \text{fermé}+\text{quasi-projectif}
\end{aligned}
\]

\diagram{\begin{tikzcd}[ampersand replacement=\&,column sep=.45em,row sep=1.05em,cells={font=\scriptsize},execute at end picture={\draw ([xshift=-.45em,yshift=.55em]\tikzcdmatrixname-1-2.north west) rectangle ([xshift=.45em,yshift=-.5em]\tikzcdmatrixname-3-7.south east);}]
 {} \&\text{imm. fermée}\arrow[rr,Rightarrow]\arrow[d,Rightarrow] \&{} \&\text{fini}\arrow[rr,Rightarrow]\arrow[dl,Rightarrow]\arrow[d,Rightarrow] \&{} \&\text{entier}\arrow[r,Rightarrow]\arrow[d,Rightarrow] \&\text{affine}\arrow[dd,Rightarrow]\\
 {} \&\text{imm. qu. cpt}\arrow[d,Rightarrow] \&\substack{\text{affine}\\\text{de t.f.}}\arrow[dl,Rightarrow] \&\text{proj}\arrow[r,Rightarrow]\arrow[d,Rightarrow] \&\text{propre}\arrow[r,Rightarrow]\arrow[drr,Rightarrow] \&\text{univ. fermé} \&{}\\
 \deleted{\substack{\text{affine}\\\text{de t.f.}}}\arrow[r,Rightarrow]\&\substack{\text{quasi-affine}\\\text{de t.f.}}\arrow[rr,Rightarrow] \&{} \&\text{qu. proj}\arrow[rrr,Rightarrow] \&{} \&{} \&\text{séparé}
\end{tikzcd}}

\uline{Th.} \textcircled{a} \uline{fini} $={}$ \uline{propre à fibres finies} (\unsure{ensemblistement|cas noethérien}) $={}$ projectif à fibres finies

\deleted{\textcircled{b} \uline{séparé} \unsure{quasi-fini|quasi-} $\Rightarrow$ \uline{\unsure{qu. affine|affine}} \unsure{à} \uline{fibres finies}}

% archive sheet number: 8
\scanpage{10}
\margin{EGA \uline{II} 4.5}

$X$ \uline{cohérent} $L$ inversible. Conditions équivalentes

a) les $X_f$ forment base de la topologie de $X$ ($f\in\unsure{S_n|\mathcal S_n}=\Gamma(X,L^{\otimes n})$, $n\geq1$)

a') les $X_f$ affines forment une base \ldots\ldots

b) Pour tt $F$ quasi-cohérent \inserted{de t.f.} sur $X$, $F(n)$ est engendré par ses sections pour $n$ grand [ou encore : $F$ \inserted{$\simeq$} un quotient d'un $\underline O_X^N(-n)$]

b') Comme b), avec $F\subset\underline O_X$

\deleted{\illegible{1} Si}

Il \uline{suffit} \deleted{que}, si $X$ est au dessus de $S$ \uline{affine}, qu'il \deleted{\unsure{existe}} \inserted{\unsure{ait|soit}}

c) $\exists n>0$ tel que $L^{\otimes n}$ très ample rel/$S$

c') $\exists n_0>0$ tel que $L^{\otimes n_0}$ très ample pour $n\geq n_0$

Cette condition est aussi nécessaire \uline{si} \deleted{$\unsure{S|X}$} $X$ \uline{de t.f. sur} $S$.

\uline{Définition} $L$ est dit ample si $L$ satisfait les conditions équivalentes a) à b'). Si \deleted{$Y$ \unsure{n'est}} on a $f:X\longrightarrow Y$, \deleted{morphisme} \deleted{cohérent} on dit que $L$ est $f$-ample (ou ample rel/$Y$) \deleted{sss} si $\exists$ recouvrement ouvert affine de $Y$ par des $Y_i$, tel que \inserted{($X|Y_i$ cohérent et)} $L|(X|Y_i)$ ample pour tt $i$.

\deleted{Il suffit} Alors c'est vrai pour tt $i$, et on peut relativiser aisément les \deleted{\unsure{aux}} critères b), b') (sur $Y$ qu. cpt du moins)

\deleted{Très ample \unsure{ou}}

\uline{NB} Si $\exists n$ tel que $L^{\otimes n}$ soit très ample rel/$Y$, alors il est ample rel/$Y$, la réciproque est vraie si $f$ est de t.f. et $Y$ qu. cpt.

\deleted{\unsure{$X$|Soit}}\uline{Prop}

\diagram{\begin{tikzcd}[ampersand replacement=\&]
 L\&M\&{}\\
 X\arrow[r,dash,"f"]\&Y\arrow[r,dash,"g"]\&Z\\
 {}\&{}\&\text{\uline{qu. cpt.}}
\end{tikzcd}}

\[
\begin{aligned}
 &\deleted{\unsure{L}}\quad L\ \text{ample rel}/Y\\
 &\quad M\ \text{ample rel}/Z
\end{aligned}
\quad\unsure{\Longrightarrow|\nRightarrow}\quad
\begin{aligned}
 &\exists n_0\ \text{tel que }n\geq n_0\\
 &\text{implique } L\otimes f^*(M^{\otimes n})\\
 &\text{ample rel}/Z.
\end{aligned}
\]

($=$ chose vaut avec très ample)

% archive sheet number: 9
\scanpage{11}
Schémas irréductibles, connexes \ldots / noethériens (attention!), loc. noeth /

réduits, intègres ($={}$ irréduct $+$ réduits) normaux \deleted{\unsure{singuliers|réguliers}} (anneaux locaux n.) \deleted{réguliers (id) \illegible{1}}

global $\left\{\begin{aligned}&\text{Morphismes qu. \uline{cpts}}\\&\text{de \uline{type fini}}\end{aligned}\right\}$ a) b) (cas de $V(M),P(M)$)

local $\left\{\begin{aligned}&\text{loc. \uline{de t.f.}}\qquad\text{a) b) c)}\\&\text{loc. \uline{de prés. fi.}}\qquad\text{a) b) }\deleted{\unsure{\text{c}|\text{c)}}}\end{aligned}\right.$

[$f$ \deleted{\unsure{de t.f.}} loc. de t.f. $\Rightarrow\operatorname{diag}_f$ loc. de prés. finie]

\uline{immersions} \qquad a) b) c)

imm. \uline{fermées}, \uline{ouvertes} [\uline{surjectives}] a) b)

sous-schémas ouverts, fermés [transitif, stable par im. réciproque]

\diagram{\begin{tikzcd}[ampersand replacement=\&]
 X\arrow[d]\&X'\arrow[l]\arrow[d,dash]\\
 Y\arrow[r,dash,"\text{cart.}" description]\&Y'
\end{tikzcd}}

$X\overset{f}{\mathrel{-}}Y\overset{g}{\mathrel{-}}Z$

\[
\begin{aligned}
 &\left\{\begin{aligned}
 &\text{a) Si } f\in M,\ f'\in M\\
 &\quad(\Rightarrow\text{ tout isom }\in M)\\
 &\text{b) Si }f,g\in M,\ gf\in M
 \end{aligned}\right.\\
 &\deleted{\unsure{\Leftarrow}}\Longrightarrow(f_1\in M,\ f_2\in M\Longrightarrow f_1\times f_2\in M)\\
 &\Longrightarrow(gf\in M,\ \operatorname{diag}(g)\in M\Longrightarrow f\in M)\quad *
\end{aligned}
\]

c) Si $gf\in M$, $f\in M$ [en présence de a) et b), équivaut à : toute section est $\in M$, \{ou : diagonales $\in M$\} ; ou : noyaux $\in\underline{M}$ si c'est local en bas \inserted{et vrai loclt sur \unsure{but}} \} on est ramené au cas affine]

\margin{Ex : vérifier sur diagonal, sections, noyaux}

morphismes \uline{séparés} [a) b) c)] : diagonale fermée $\Longleftrightarrow$ noyaux de couples dedans fermés $\Longleftrightarrow$ principe d'identité pour morphismes des schémas \uline{réduits} dedans.

\uline{Ex} (mono -- donc applic. diag, noyaux, diag. morphismes affines

\deleted{\unsure{morphismes quasi-séparés}}

$P(V)\longrightarrow S$

Schémas séparés (ils sont séparés sur n'importe quoi). $\Rightarrow$ \uline{Intersection} \deleted{\unsure{affines}} de deux ouverts affines affine sur $S$ i.e. $X\unsure{,}Y$ sur $S$-affine $\Rightarrow X\times_S Y$ affine

% slanted
\margin{on \uline{peut} \uline{relativiser} $gf$ affine, $g$ séparé $\Rightarrow f$ affine}

morphismes \uline{quasi-séparés} (a) b) c)) (\deleted{\unsure{not}} schémas quasi-séparés)

diagonale quasi-cpte $\Longleftrightarrow$ noyaux de couples dedans quasi-compacts $\Longleftrightarrow$ sections après chgt de base qu. cpt

\deleted{\unsure{et la base}} [Si base qu. sép : $\Longleftrightarrow$ \deleted{\unsure{la diagonale qu. cpt}} l'intersection de 2 ouverts qu. cpts est \unsure{enc. q|qu. cpte}

\deleted{\unsure{morphismes} \unsure{séparés} (\unsure{cf.}) \unsure{et qu. ctp et quasi-séparés})}

$\Longleftrightarrow$ \deleted{tout} schém. \unsure{sur} $X$ qu. cpt. est qu. compact sur $X$.]

% slanted
\margin{$gf$ qu. cpt (de t.f.) $g$ qu. sép $\Rightarrow$ $f$ qu. cpt (qu. sép)}

morphismes \uline{cohérents} $={}$ qu. cpt qu. sép \quad a) b) [$gf$ coh et $g$ qu. sép $\Rightarrow f$ coh]

Si $f:X\mathrel{-}Y$ qu. coh. et $F$ qu. coh. sur $X$, alors $f_*(F)$ (et les $R^i f_*(F)$) quasi-cohérent.

\margin{schémas cohérents}

\uline{NB} Tous les schémas envisagés par les old-timers étaient de t.f. et \uline{séparés} sur des corps. La plupart des schémas qu'on envisage \unsure{encore aujourd'hui|encore auj.|couramment} sont séparés pratiquement \uline{tous} ceux qu'on rencontre sont quasi-séparés.

Beaucoup sont qu. cpts -- mais il y en a d'importants qui ne le sont pas.

[morphismes de \uline{prés. finie} ($={}$ loc. de prés. finie et cohérents) \quad a) b) \quad [$gf$ de prés. finie et $g$ de t.f et qu. sép.

\fbox{Toutes les notions sont \uline{locales en bas}}

$\longrightarrow$ Yoga des passages à la limite.

\uline{Généralités} \quad Local en bas $\left\{\begin{aligned}&\text{stable par image inverse}\\&\text{stable par composition}\end{aligned}\right.$ $\substack{\Longrightarrow\\\Longrightarrow}$

$g:Y\mathrel{-}S$

\[
\begin{aligned}
\operatorname{diag}g\in M &\Longleftrightarrow \forall S'\mathrel{-}S,\ \text{tte section de }Y'\text{ sur }S'\text{ est }\in M\\
 &\Longleftrightarrow \forall (u,v):X\rightrightarrows Y\text{ sur }S,\quad\ker(u,v)\lhook\joinrel\longrightarrow X\text{ est }\in M\\
 &\Longleftrightarrow \forall f:X\mathrel{-}Y,\quad\Gamma_f:X\mathrel{-}X\times_S Y\text{ est }\in M\\
 &\Longleftrightarrow gf\in M\Longrightarrow f\in M
\end{aligned}
\]

\diagram{\begin{tikzcd}[ampersand replacement=\&,column sep=small]
 *\quad X\arrow[d,"f"']\&{}\\
 Y\arrow[d,"g"']\&X\times_S Y\arrow[l,"\mathrm{pr}_2"']\arrow[d,"\mathrm{pr}_1"]\\
 S\&X\arrow[l,"gf"]\arrow[u,bend right,"\Gamma_f"']
\end{tikzcd}}

\diagram{\begin{tikzcd}[ampersand replacement=\&]
 X\arrow[r,"\Gamma_f"]\arrow[d,"f"']\&X\times_S Y\arrow[d,"f\times\mathrm{id}_Y"]\\
 Y\arrow[r,hook,"\operatorname{diag}_Y"']\&Y\times_S Y
\end{tikzcd}}

\uline{plan} \quad a) Traversée du désert

b) $\left\{\begin{aligned}&\text{Passage au quotient}\\&\text{descente fid. plat}\end{aligned}\right\}$ et langage faisceautique

c) Notions ``géométriques'' sur un corps

d) \deleted{\illegible{1}} Ensembles constructibles, passages à la limite

% archive sheet number: 10
\scanpage{12}
\deleted{Localisé d'un schéma en un point, schéma local, morphisme d'un schéma local dans un schéma.}

\deleted{I Morphismes quasi-compacts de type fini loc. de type fini loc. de prés. finie}

\diagram{\begin{tikzcd}[ampersand replacement=\&]
 \deleted{X}\arrow[d,"f"']\&\deleted{X'}\arrow[l]\arrow[d,"f'"]\\
 \deleted{Y}\arrow[r,dash]\&\deleted{Y'}
\end{tikzcd}}

\diagram{\begin{tikzcd}[ampersand replacement=\&]
 \deleted{X}\arrow[rr,"f"]\arrow[dr,dash]\&{}\&\deleted{Y}\arrow[dl,dash,"g"]\\
 {}\&\deleted{S}\&{}
\end{tikzcd}}

\deleted{Si $f$ est $M$, $f'$ est $M$. Si $f$ et $g$ sont $M$, $fg$ aussi $\Rightarrow$ Si $f_1$ et $f_2$ sont $M$, $f_1\times_S f_2$ est $M$}

\deleted{Si $gf$ est $M$, $f$ est $M$] vrai pour \unsure{loc. de t.f.|loc. t.f.}, \deleted{\unsure{loc. prés.}} imm., morphismes séparés}

\deleted{immersions $\left\{\begin{aligned}&\text{ouvertes}\\&\text{fermées}\end{aligned}\right.$}

\margin{\deleted{immersion surjective \illegible{4} \unsure{ou} \illegible{2}}}

\deleted{Sous-schémas ouverts -- fermés}

\deleted{\uline{Ex} \{ Morphisme diagonal -- noyaux \quad sections}

\deleted{morphismes séparés : diagonale fermée -- noyaux fermés -- principe de prolongement des identités $Z\longrightarrow X$ pour $Z$ réduit ($X$ séparé sur $S$)}

\deleted{morphismes qu}

\textcircled{II} $\left\{\begin{aligned}
 &\text{Schémas noeth, loc. noeth, artiniens, loc. artiniens}\\
 &\text{Schémas irréd, connexes}\\
 &\text{Schémas \deleted{réduits}, intègres}\\
 &\text{Schémas normaux, réguliers}
\end{aligned}\right.$

I) \unsure{Schéma localisé|Schémas localisé|Schémas localisés} d'un \deleted{\unsure{sc}} schéma en un point schémas locaux ($\exists!$ pt fermé, et quasi-compact) Morphisme d'un schéma \inserted{local} dans un schéma.

\textcircled{III} Morphismes : \uline{affines}, \inserted{[mono, surjectifs]} [\uline{qu. cpt}, \uline{de t.f}, \uline{loc.} de t.f., \uline{loc.} de \uline{p.f.}], \uline{imm.} (\uline{ouv.}, \uline{fermées}, \uline{surjectives}), \uline{séparés}, [\uline{quasi-séparés} -- \uline{cohérents}, de \uline{prés. finie}]

\textcircled{c} \uline{morphismes} \inserted{[\uline{entiers}]} \uline{finis}, [quasi-finis]

\deleted{\unsure{quasi}} \textcircled{a} (\uline{projectifs}, \uline{quasi-projectifs}) [quasi-affines]

\textcircled{d} \uline{radiciels} homéom. universels \deleted{\unsure{univ|univ.}} ouverts, univ. fermés

\textcircled{b} \uline{propres}

\margin{Sur base cohérente}

\uline{Oubli} : \{ Cohérence du faisceau structural sur un schéma loc. noeth, \uline{structure induite réduite} sur un \uline{fermé}

\inserted{attention, une imm. ouverte n'est pas une imm. \unsure{quasi entière|quasi-entière}, \unsure{idée} \unsure{mettre} : cas des deux \unsure{variétés même corps|valuations même corps})}

\inserted{où on rencontre les schémas non séparés}

Un schéma loc. noeth. est qu. sép.

$f$ cohérent, $F$ qu. coh $\Rightarrow f_*(F)$ qu. coh.

\uline{cor} \{ Extension de faisceaux qu. coh sur ouvert rétrocompact

Extension de sous-faisceaux \deleted{\unsure{de t.f. de type}} \unsure{de} $F|U$ ($U$ rétro-cpt)

$\left\{\begin{aligned}
 &\text{\uline{net}, \uline{lisse}, \uline{étale}}\\
 &\boxed{\text{\uline{plat}, \uline{fid plat}}}
\end{aligned}\right.$

Propriétés géom des fibres \ldots

% archive sheet number: 11
\scanpage{13}
\deleted{\unsure{Morphismes quasi-projectifs}}

\unsure{Épim. couvrants|Épimorphismes|épimorphismes universels} (Zar) de schémas (= foncteurs représentables)

% slanted
\margin{\textcircled{b} faisceaux sur (Sch)}

\textcircled{d} Schém. de type fini sur un autre \inserted{$S$} : [quasi-comp. et sur $S$ et loc. de type fini sur $S$]

\textcircled{a} Schémas comme foncteurs : avantages respectifs des arguments anneaux (schémas affines) et des arguments schémas (\uline{voire espaces localement annelés}).

\textcircled{e} Corollaire du th. de représentabilité : nature \unsure{localement|locale} du \unsure{pb.|Pb.}

\textcircled{f} Schéma $P(M)$, $M$ module qu.-coh sur $S$ \inserted{(pas néc. affine)} : il représente un foncteur remarquable \unsure{m.} dans la catégorie des espaces loc. annelés sur $S$ (voire des topos loc. annelés sur $S$). Sa structure essentielle : --

% slanted
\margin{c'est un th. de représentabilité}

Dépendance p.r. au changement de base ; dépendance fonctorielle en $S$ / Dépendance en $M$ (pour épi)

[Schémas \inserted{quasi-}projectifs sur $S$ ($S$ affine disons) : qui se plongent \deleted{\unsure{d}} dans un $P(M)$, \deleted{\unsure{de de de présentation}} \inserted{(on peut alors choisir $V$ de type fini)}]

\inserted{de t.f. sur $S$,}

Relation avec $V(M)$ (fibré vectoriel sur $S$ associé à $M$)

\[ V(M)^*\subset V(M)\qquad(\text{les \uline{épi}}\quad M_{k'}\longrightarrow k') \]

\[ V(M)^*\longrightarrow P(M)\qquad\text{\uline{épi de faisceaux}}\quad[\text{\uline{i.e. il y a une section loc pour Zar.}}] \]

En fait, le $\unsure{\mathrm{Gr}_m|G_m|\mathbb{G}_{m,S}}$ opère sur $V(M)^{\deleted{*}}$, laissant $V(M)^*$ stable ; il y opère librement, et

\[ P(M)\simeq V(M)^*/G_m\quad(\text{isom de faisceaux Zar}) \]

\textcircled{g} $\varprojlim$ finies dans (Sch)

\textcircled{h} Spectre d'Algèbre qu. coh. et \deleted{\unsure{schéma}} morphismes affines (les schémas affines sur $S$)

% archive sheet number: 12
\scanpage{15}
1) Morphisme d'un espace loc annelé dans schéma affine

\[
 \operatorname{Hom}_k(X,\underset{\operatorname{Spec} A}{\underset{\shortparallel}{Y}})
 \mathrel{\substack{\overset{\sim}{\longrightarrow}\\[-2pt]\longleftarrow}}
 \operatorname{Hom}_{k\text{-alg}}(\underset{\Gamma(Y,\underline O_Y)}{\underset{\shortparallel}{A}},\Gamma(X,\underline O_{\deleted{\unsure{x}}X}))
\]

\deleted{\(\forall x\in\)} Soit $u:A\mathrel{-}\Gamma(X,\underline O_X)$

Pour tout $x\in X$, soit $\varphi(x)$ \inserted{$\in Y=\operatorname{Spec}(A)$} le noyau de

\diagram{\begin{tikzcd}[ampersand replacement=\&]
 A\arrow[r,dash,"u"]\&\Gamma(X,\underline O_{\unsure{X|x|\mathcal X}})\arrow[rr,dash,"\varepsilon_x"]\arrow[dr]\&{}\&k(x)\\
 {}\&{}\&\underline O_{X,x}\arrow[ur]\&{}
\end{tikzcd}}

d'où $\varphi:X\longrightarrow Y$

Si $f\in A$, $\varphi^{-1}(U_f)={}$ ens des $x\in X$ tels que \inserted{$\varepsilon_x$}$(u(f))\ne 0$ i.e. tels que $u(f)_x$ inversible, i.e. $X_{u(f)}$ : il est ouvert donc $\varphi$ continu.

De plus, on trouve

\[ A_{\deleted{\unsure{g}}f}\longrightarrow\Gamma(U_f,\underline O_X)\qquad\text{rend $f$ inversible, \unsure{donc|donc se|en} \unsure{se factorise en|factorise en}} \]

\[ \Gamma(Y_f,\underline O_Y)\simeq A_f\mathrel{-}\Gamma(U_f,\underline O_X)\quad\text{avec compatibilité pour \unsure{restriction|restrictions}, d'où un morphisme annelé.} \]

Il est \uline{loc}. \uline{annelé}, car \deleted{\unsure{par définition, il transforme}} \deleted{\unsure{un}} on a $\varphi^{\unsure{-1|*}}(Y_f)=X_{u(f)}$ et de même après restriction à un ouvert \ldots

\uline{Conséquence} La catégorie $(\mathrm{Sch})_{/k}$ se réconstitue connaissant $\mathrm{Sch}=\mathrm{Sch}_{/\unsure{Z|\mathbb Z}}$ - et $S=\operatorname{Spec} k$ comme $(\mathrm{Sch})_{/S}$ : les $k$-schémas sont les $S$-schémas. Mais $\mathrm{Sch}_{/S}$ important sans que $S$ soit néc. affine.

\diagram{\begin{tikzcd}[ampersand replacement=\&,column sep=small,row sep=small]
 (\mathrm{Sch}\text{-Aff})_k^{\unsure{o|\circ}}\arrow[d,phantom,"\shortparallel"]\&{}\&{}\\
 (\mathrm{Aff}_k)\arrow[r,hook,"\mathrm{incl}"]\arrow[d,phantom,"\shortparallel"]\&(\mathrm{Sch})_{/k}\arrow[r,dash,hook']\&(\mathrm{Locan})_k\arrow[dl]\\
 (\mathrm{Alg}_k)^o\&\widetilde{(\mathrm{Sch})_{/k}}\arrow[r,hook,"\mathrm{restr.}","{\substack{\text{faisceaux}\\\text{pour ``top.}\\\text{de Zariski''}}}"']\&\widetilde{\mathrm{Aff}}_k\simeq\deleted{\unsure{\widetilde{\mathrm{Alg}}_k}}\underline{\operatorname{Hom}}^{\sim}(\mathrm{Alg}_k,\mathrm{Ens})\\
 \widetilde{(\mathrm{Aff}_k)}\arrow[uu,hook]\arrow[d,phantom,"\shortparallel"]\&{}\&{}\\
 \underline{\operatorname{Hom}}^{\sim}(\mathrm{Alg}_k,\mathrm{Ens})\&\widehat{(\mathrm{Sch})_{/k}}\arrow[uu,hook,"\mathrm{inc}"]\arrow[r,"\mathrm{restr.}","(\text{\unsure{pas pl. fid|pas pl. fid.|pas plein fid.}})"']\&\widehat{\mathrm{Aff}}_k\simeq\deleted{\unsure{\mathrm{Alg}|\mathrm{Alg}_k}}\underline{\operatorname{Hom}}(\mathrm{Alg}_k,\mathrm{Ens})\arrow[uu,hook,"\mathrm{inc}"]\\
 \widehat{(\mathrm{Aff}_k)}\arrow[uu,hook,"\mathrm{inc}."]\arrow[d,phantom,"\shortparallel"]\&{}\&{}\\
 \underline{\operatorname{Hom}}(\mathrm{Alg}_k,\mathrm{Ens})\&{}\&{}
\end{tikzcd}}

Comment exprimer en termes d'algèbre commutative qu'un \inserted{co-}foncteur sur $\mathrm{Aff}_k$ est un faisceau, ? i.e. un foncteur Aff

[ Interprétation en termes de foncteurs d'anneaux de

\diagram{\begin{tikzcd}[ampersand replacement=\&]
 {}\&(\mathrm{Sch})_k\&{\simeq}\&\mathrm{Sch}_{/\operatorname{Spec} k=S}\\
 C=\mathrm{Aff}_k\&\widehat{C_{/S}}\arrow[u,hook]\&{\simeq}\&\widehat C_{/\underline S}\arrow[u,hook]
\end{tikzcd}}
]

$\left\{\begin{aligned}&\text{Interprétation des \uline{sommes} en termes faisceautiques}\\
 &\text{(pas d'ennui pour les \deleted{produits} $\varprojlim$ (finies, \inserted{disons} \ldots.)}
 \end{aligned}\right.$

\margin{Notion de Sup de sous-faisceaux}

% archive sheet number: 14
\scanpage{16}
\deleted{Or on vérifie facilement \uline{directement} \inserted{(le faire ....)} que (pour les sous-catégories str. pleines envisagées \unsure{en|de} $\unsure{\pi|\Pi}(C',\widehat D)$ du diagramme}

\deleted{$\pi(C,\widetilde D)=$}

\diagram{\begin{tikzcd}[ampersand replacement=\&]
 {}\&\deleted{\pi(C,\widetilde D)}\arrow[dl,dash]\arrow[dr]\&{}\\
 \deleted{\pi(C',\widetilde D)}\arrow[dr,dash]\&{}\&\deleted{\pi(C,\widehat D)}\arrow[dl,dash]\\
 {}\&\deleted{\pi(C',\widehat D)}\&{}
\end{tikzcd}}

\deleted{on a}

\[\deleted{\pi(C,\widetilde D)=\pi(C',\widetilde D)\cap\pi(C,\widehat D)}\]

\deleted{donc on a ramené, grâce à ? \unsure{relation|obligation|utilisation} analogue}

\[\deleted{C\underset{\mathrm{Top}}{\overset{\deleted{\unsure{\times}}}{\times}}\widetilde D=(C'\underset{\mathrm{Top}}{\times}\widetilde D)\cap(C\underset{\mathrm{Top}}{\times}\widehat D)}\]

\deleted{à prouver que $\gamma(C,\widehat D)$ et $\gamma(C',\widetilde D)$ sont des équivalences i.e. \unsure{les|des} foncteurs $\gamma(C,\widehat D)_{\unsure{*|K}}$ et $\gamma(C',\widetilde D)_*$ ess. surj. ou encore}

% archive sheet number: 15
\scanpage{17}
Caractérisation des foncteurs $(\mathrm{Sch})_{/k}^{o}\longrightarrow\mathrm{Ens}$ qui sont repr

\unsure{Idem|Id.} : \rule{1.7cm}{0.2pt}\quad $\mathrm{Alg}_k\mathrel{-}\mathrm{Ens}$ \rule{2cm}{0.2pt} ].

a) Immersions ouvertes de faisceaux

[stable par chgt de base, par composition]

b) Sup de sous-faisceaux d'un faisceau.

\uline{Prop.} Soit $F$ un foncteur $(\mathrm{Sch})_{/S}\mathrel{-}(\mathrm{Ens})$. Pour que $F$ soit repr. \inserted{(i.e. un \unsure{schéma|schém.})}, il f et s. qu'il satisfasse les conditions suivantes

a) $F$ un faisceau

b) $\exists\ (F_i)_{i\in I}$ famille de \uline{sous faisceaux} \uline{ouverts}, tel que $F=\operatorname{Sup} F_i$ \inserted{(ds cat. de faisceaux)}, et chacun des $F_i$ repr. \unsure{(i.e. un schéma)|par un schéma|i.e. un schéma}

\uline{Cor} On peut \deleted{\unsure{donner|énoncer}} remplacer b) par

b') comme b), les $F_i$ étant de plus des schémas \uline{affines}.

\uline{Cor} (Supposant $S=\operatorname{Spec} k$) Pour qu'un foncteur $k\text{-}\mathrm{Alg}\longrightarrow(\mathrm{Ens})$ soit un faisceau il f et s. que \deleted{\unsure{il soit}} \inserted{ce soit} un \deleted{\unsure{préfais}} \deleted{\unsure{ceau}} faisceau, et qu'on le puisse recouvrir par des sous-faisceaux ouverts affines ....

Ainsi, la notion de schéma est explicitée sans recours à la notion d'espaces topologiques, faisceau d'anneaux, spectre d'un anneau ...

\uline{Applications} a) \uline{Représentabilité du foncteur} $P(E)$

($E$ un module loc. libre sur $S$ -- p. ex. $S=\operatorname{Spec} k$, $E$ un module de \unsure{r.f.|t.f.} sur $k$)

b) \uline{Représentabilité} de \uline{Spec} $A$

c) \uline{Produits fibrés} (donc $\varprojlim$ finies)

% slanted
\margin{sommes sont disjointes et universelles}

d) $\varprojlim$ \inserted{quel.} pour des schémas affines au dessus d'un schéma

% archive sheet number: 16
